%%%%

%%%   Самарский государственный технический университет

%%%   Кафедра прикладной математики и информатики

%%%

%%%   Преамбула для статьи в журнал "Вестник Самарского государственного

%%%   технического университета. Серия: «Физико-математические науки»"

%%%   Ориентирована под MikTEx 2.4 for Win32

%%%

%%%   Хотя сам журнал собирается под GNU/Linux на дистрибутиве TeXLive



\documentclass[11pt,a4paper,twoside]{article}



\usepackage[cp1251]{inputenc}

\usepackage[T2A]{fontenc}

\usepackage[english,russian]{babel}

\usepackage{samgtubib}

\usepackage[dvips]{graphicx}

\usepackage[multiple]{footmisc}



\usepackage{samgtu}

\renewcommand{\VestnikYear}{2009} % Год издания Вестника

\renewcommand{\VestnikNumber}{2\,(19)} % Номер Вестника

\renewcommand{\VestnikMonth}{Ноябрь} % Месяц выхода Вестника

\renewcommand{\VestnikData}{01.11.09} % Дата подписания Вестника в печать

\renewcommand{\VestnikUPL}{26,64} % Усл. печ. л.

\renewcommand{\VestnikUIL}{25,73} % Уч.-изд. л.

\renewcommand{\VestnikRegNumber}{479/09} % Регистрационный номер

\renewcommand{\VestnikOrder}{994} % Номер заказа






%
%
%\begin{document}


\renewcommand{\firstpage}{194}
\renewcommand{\lastpage}{199}

\udk{517.51+517.98}
\msc{42C15}

\title{О представлении фреймов Парсеваля}{О представлении фреймов Парсеваля}

\author{И.\,С.~Рябцов}{Р\,я\,б\,ц\,о\,в~И.\,С.}

\institute{\samgu.}

\email{E-mails}{tinnulion@mail.ru}

\otherinf{\fullauthor{Игорь Сергеевич Рябцов}  \position{аспирант} \dept{каф. функционального анализа и~теории функций}}

\keywords{фреймы Парсеваля, эквивалентность фреймов, представление фреймов, равноугольные фреймы, жёсткие фреймы}

\workabstract{Работа посвящена исследованию свойств фреймов Парсеваля в~конечномерных пространствах, а~именно возможности
представления одних фреймов как суммы других. Даётся новый подход к построению произвольных фреймов Парсеваля, а также
описывается алгоритм разложения произвольного фрейма в~сумму. В~работе описывается ряд особых свойств равноугольных жёстких
фреймов применительно к~поставленным задачам.}

\engtitle{On representation of Parseval frames}

\engauthor{I.\,S.~Ryabtsov}{R\,y\,a\,b\,t\,s\,o\,v~I.\,S.}

\enginstitute{\engsamgu.}

\otherenginf{\fullauthor{Igor S. Ryabtsov} \position{Postgraduate Student} \dept{Dept. of Functional Analysis and 
Fucntions Theory}}

\engabstract{This paper investigates properties of Parseval frames in finite dimensional vector spaces, namely, the possibility of representing some frames as sums of others. A new approach in constructing arbitrary Parseval frames and the decomposition arbitrary frame into the sum are described. Besides there is a number of special properties of equiangular tight frames.}

\engkeywords{Parseval frames, frame equivalency, frame representations, equiangular frames, tight frames}

\date{20/XII/2010}
\revisiondate{10/V/2011}

\maketitle

\Section[n]{Введение} 
Пусть $M$, $N$ "--- натуральные числа, $M \geq N.$ 
Пространство $\mathbb{R}^N$ наделяется стандартным
скалярным произведением $\langle {}\cdot{}, {}\cdot{} \rangle$ и согласованной нормой $\|x\|=\sqrt{\langle x,\,x\rangle}.$
\begin{definition}
Набор элементов $F =\{f_i \}_{i=1}^{M}$ называется {\it фреймом} в пространстве $\mathbb{R}^N$, если существуют числа $A$, $B>0$
такие, что для всех $x \in \mathbb{R}^N$ выполняется двойное неравенство
$$
A \|x \|^{2} \leq \sum_{i=1}^M |\langle x,\,f_i \rangle |^{2} \leq B\|x\|^{2}.
$$
\end{definition}
Оговоримся, что среди векторов фрейма нет нулевых и коллинеарных векторов. В данной
работе ограничимся рассмотрением фреймов в конечномерных пространствах с конечным числом элементов.

Числа $A$ и $B$ называются границами фрейма. Они определены неоднозначно, cупремум множества всех нижних границ и инфимум
множества всех верхних границ фрейма называются {\it оптимальными границами фрейма}. Если оптимальные нижние и верхние границы
фрейма совпадают и выполняется равенство
\begin{equation}
\label{ryab:eq1}
A \|x \|^{2} = \sum_{i=1}^M | \langle  x,\,f_i \rangle |^{2},
\end{equation}
то будем называть такой фрейм {\it жёстким}. Если при этом $A = 1$, то фрейм будем называть {\it фреймом Парсеваля}. Для жёстких фреймов выполняется
равенство, которое эквивалентно равенству \eqref{ryab:eq1}:
\begin{align}
\label{ryab:eq11}
Ax = \sum_{i=1}^M \langle x,\, f_i \rangle f_i.
\tag{$1'$}
\end{align}
Если норма каждого вектора фрейма равна единице, то фрейм будем называть {\it нормированным}.
Доказательство приведенных утверждений и более полную информацию по теории фреймов можно
найти в работах \cite{ryabtsov:Cri,ryabtsov:Cas}.

\begin{definition}
Определим оператор $T$ для произвольного фрейма $F = \{f_i \}_{i=1}^M$. Пусть существует пара векторов
$f_k$, $f_p \in F$ с $\langle f_k,\, f_p \rangle = \|f_k \| \|f_p\|$
и $k \neq p$. Тогда оператор действует на фрейм по следующему правилу:
\begin{equation}
\label{ryab:eq2}
T (F ) = F \backslash \{f_k,\,f_p \} \cup \left\{ \sqrt{ \|f_k \|^2 + \|f_p \|^2} \frac{f_k}{\|f_k \|} \right\}.
\end{equation}
Если такой пары векторов не существует, то $T (F ) = F$. Действие оператора на фрейм можно упростить по формуле
\begin{equation}
T (F ) = F \backslash \{f_k,\, f_p \} \cup  T \left( \{f_k,\, f_p \}  \right).
\tag{$2'$}
\end{equation}

Легко видеть, что этот оператор фактически заменяет пару коллинеарных векторов на один.
\end{definition}

\begin{definition}
Поскольку число коллинеарных векторов в конечном фрейме конечно, то, применяя оператор $T$ достаточное число раз,
получим фрейм без коллинеарных векторов. Обозначим такой оператор $T_\infty$:
$$
T_\infty = T \circ \cdots \circ T .
$$
\end{definition}

\smallskip

\Section[n]{Основные результаты}

\begin{definition}
Будем называть фрейм Парсеваля $F = \{f_i\}_{i=1}^M$ {\it составным}, если существует набор неотрицательных констант
$\{\alpha_i\}_{i=1}^M$ такой, что система векторов
$F_\alpha = \{\alpha_i f_i\}_{i=1}^M$ "--- также фрейм Парсеваля, при этом хотя бы одно $\alpha_i$ равно нулю.
Будем называть фрейм Парсеваля $F = \{f_i\}_{i=1}^M$ {\it простым}, если он не является составным.
\end{definition}

\begin{theorem}
Для любого простого фрейма Парсеваля $F = \{f_i\}_{i=1}^M$ и ортогональной матрицы $Q$
система векторов $\{Q f_i \}_{i=1}^M$ является простым фреймом Парсеваля.
\end{theorem}

\begin{proof}
Проведём доказательство методом от противного. Предположим, что $\{f_i\}_{i=1}^M$ "--- простой фрейм Парсеваля, а $\{Q f_i\}_{i=1}^M$
"--- составной фрейм. Тогда по определению существует набор коэффициентов трансформации $\{\alpha_i\}_{i=1}^M$ такой, что
фрейм $\{\alpha_i Q f_i\}_{i=1}^M$ "--- фрейм Парсеваля, при этом какая"=то из констант $\alpha_i$ равна нулю. Но тогда согласно
определению фрейма Парсеваля имеем
\begin{multline*}
x = \sum_{i=1}^M  \langle  x,\, \alpha_i Q f_i \rangle \alpha_i Q f_i =
Q \sum_{i=1}^M \langle  Q Q^{\top} x,\, Q \alpha_i f_i \rangle \alpha_i f_i =\\ =
Q \sum_{i=1}^M \langle  Q^{\top} x,\, \alpha_i f_i \rangle \alpha_i f_i =
\sum_{i=1}^M \langle  x,\, \alpha_i f_i \rangle \alpha_i f_i.
\end{multline*}
Получаем, что система $\{f_i \}_{i=1}^M$ "--- составной фрейм Парсеваля, но это противоречит первоначальному предположению.
\end{proof}

\begin{definition}
{\it Суммой} пары фреймов Парсеваля $F_1 = \left\{f_i^{F_1}\right\}_{i=1}^{M_1}$ и $F_2 = \left\{f_i^{F_2}\right\}_{i=1}^{M_2}$ будем называть
фрейм Парсеваля, который состоит из векторов
$$
F_1 \oplus F_2 = T_\infty \left(\left\{ f_i^{F_1} \right\}_{i=1}^{M_1} \cup \left\{f_i^{F_2}\right\}_{i=1}^{M_2} \right).
$$

Число векторов в результирующем фрейме не превосходит $M_1 + M_2$.
Данная операция является коммутативной, ассоциативной.
\end{definition}

Аналогично можно определить сумму конечного числа фреймов c аналогичными свойствами:
$$
\overset{K}{\underset{k = 1}{\oplus}} F_k = T_\infty\left(\overset{K}{\underset{k = 1}{\cup}} \left\{
f_i^{F_k}\right\}_{i=1}^{M_k}\right), \quad
F_k = \overset{K}{\underset{k = 1}{\cup}}  \left\{f_i^{F_k}\right\}_{i=1}^{M_k}.
$$
Взвешенная сумма $(\lambda_1 F_1)  \oplus  (\lambda_2 F_2 )$ пары фреймов Парсеваля $F_1$ и $F_2$
при условии $\lambda_1^2 + \lambda_2^2 = 1$ снова даёт фрейм Парсеваля.
\begin{theorem}
Любой составной фрейм Парсеваля $F = \{f_i \}_{i=1}^M$ можно представить как сумму конечного числа простых фреймов Парсеваля.
\end{theorem}

\begin{proof}
Для начала докажем, что любой составной фрейм Парсеваля можно представить в виде суммы двух простых или сложных фреймов.
Действительно, по определению составного фрейма существует набор констант $\{\alpha_i \}_{i=1}^M$ такой,
что $\exists k : 1 \leq k \leq M$, для которого $\alpha_k = 0$, и система векторов $\left\{\alpha_i f_i\right\}_{i=1}^M$
является фреймом Парсеваля.

Рассмотрим {\it двойственный} набор коэффициентов $\{\beta_i\}_{i=1}^M$:
$$
\beta_i = \sqrt{\frac{\alpha_{\max}^2 - \alpha_i^2}{\alpha_{\max}^2 - 1}}, \quad \alpha_{\max} = \max_{1 \leq i \leq M}{|\alpha_i|}.
$$
Двойственная система $F_\beta = \left\{\beta_i f_i\right\}_{i=1}^M$ является фреймом Парсеваля:
\begin{multline*}
\sum_{i=1}^M \langle  x,\, \beta_i f_i \rangle \beta_i f_i
= \sum_{i=1}^M \langle x,\, \sqrt{\frac{\alpha_{\max}^2 - \alpha_i^2}{\alpha_{\max}^2 - 1}} f_i  \rangle \sqrt{\frac{\alpha_{\max}^2
- \alpha_i^2}{\alpha_{\max}^2 - 1}} f_i =\\
= \frac{1}{\alpha_{\max}^2 - 1} \sum_{i=1}^M \left(\alpha_{\max}^2 - \alpha_i^2\right) \langle  x,\, f_i \rangle f_i = \\
= \frac{1}{\alpha_{\max}^2 - 1} \left( \alpha_{\max}^2 \sum_{i=1}^M \langle x,\, f_i \rangle f_i -
\sum_{i=1}^M \langle x,\, \alpha_i f_i \rangle \alpha_i f_i \right) = \\
= \frac{1}{\alpha_{\max}^2 - 1} \left( \alpha_{\max}^2 x - x \right) = x.
\end{multline*}

Покажем корректность определения набора $\{\beta_i\}_{i=1}^M$. Докажем, что $\forall k: 1 \leq k \leq M$ константа $\beta_k$ имеет
смысл и является вещественным числом, при этом набор $\{\alpha_i\}_{i=1}^M$ не совпадает с набором $\{\beta_i\}_{i=1}^M$.
Для этого достаточно доказать, что $\alpha_{\max} > 1$ при любом выборе системы $\{\alpha_i\}_{i=1}^M$, так как числитель
в силу своей природы неотрицателен.

Предположим противное, что существует набор констант $\{\alpha_i\}_{i=1}^M$ такой, что $\exists k: 1 \leq k \leq M$, для которого
$\alpha_k = 0$, и система векторов $\{\alpha_i f_i\}_{i=1}^M$ является фреймом Парсеваля, но при этом $\alpha_{\max} \leq 1$.
Оценим сверху границу фрейма $\{\alpha_i f_i\}_{i=1}^M$:
$$
\sum_{i=1}^M \left| \langle  x,\, \alpha_i f_i \rangle \right|^2 = \sum_{i=1}^M | \alpha_i |^2 \left| \langle  x,\, f_i \rangle \right|^2
< \alpha_{\max}^2 \sum_{i=1}^M \left| \langle  x,\, f_i \rangle \right|^2 = \alpha_{\max}^2.
$$
Строгое неравенство обеспечивается наличием нулевых компонент в наборе $\{\alpha_i\}_{i=1}^M$.
Противоречие с первоначальным предположением доказывает то, что $\alpha_{\max} > 1$.

Для любой пары двойственных наборов $\{\alpha_i\}_{i=1}^M$ и $\{\beta_i\}_{i=1}^M$ существуют номера
$1 \leq k \leq M$ и $1 \leq p \leq M$ такие, что $k \neq p$ и при этом
$$
\alpha_k = 0, \quad \beta_k \neq 0,  \quad \alpha_p \neq 0, \quad \beta_p = 0.
$$
Это свойство следует из определения соотвествующих наборов. Оно завершает доказательство корректности определения набора
$\{\beta_i\}_{i=1}^M$.

Любой составной фрейм Парсеваля раскладывается в объединение фреймов $\{\alpha_i f_i\}_{i=1}^M$ и
$\{\beta_i f_i\}_{i=1}^M$. Возьмём конкретные значения $\lambda_{\alpha}$ и $\lambda_{\beta}$:
$$
\lambda_{\alpha} = \frac{1}{\alpha_{\max}}, \quad
\lambda_{\beta} = \frac{\sqrt{\alpha_{\max}^2 - 1}}{\alpha_{\max}}, \quad
\lambda_{\alpha}^2 + \lambda_{\beta}^2 = \frac{1}{\alpha_{\max}^2} + \frac{\alpha_{\max}^2 - 1}{\alpha_{\max}^2} = 1
$$
и вычислим $F_\alpha \oplus F_\beta$. Поскольку фрейм $F$ коллинеарных векторов не содержит, то единственными
парами коллинеарных векторов в сумме могут быть только $\alpha_k f_k$ и $\beta_k f_k$ для $1 \leq k \leq M$.
Подействуем на эту пару оператором $T$ согласно \eqref{ryab:eq2}:
\begin{gather*}
T \left(\{ \lambda_{\alpha} \alpha_i f_i,\, \lambda_{\beta} \beta_i f_i \}\right) =
\sqrt{\|\lambda_{\alpha} \alpha_i f_i\|^2 + \|\lambda_{\beta} \beta_i f_i\|^2}
\frac{\lambda_{\alpha} \alpha_i f_i}{\|\lambda_{\alpha} \alpha_i f_i\|}
= \sqrt{\lambda_{\alpha}^2 \alpha_i^2  + \lambda_{\beta}^2 \beta_i^2} f_i,\\
\lambda_{\alpha}^2 \alpha_i^2  + \lambda_{\beta}^2 \beta_i^2 =
\frac{\alpha_i^2}{\alpha_{\max}^2} +
\frac{\alpha_{\max}^2 - 1}{\alpha_{\max}^2} \cdot \frac{\alpha_{\max}^2 - \alpha_i^2}{\alpha_{\max}^2 - 1} =
\frac{\alpha_i^2 + \alpha_{\max}^2 - \alpha_i^2}{\alpha_{\max}^2} = 1.
\end{gather*}

Таким образом, получим представление фрейма $F$ в виде
$$
F =(\lambda_{\alpha} F_\alpha) \oplus \left(\lambda_{\beta} F_\beta\right).
$$
Если полученные фреймы $F_\alpha$ и $F_\beta$ являются простыми, то процесс разложения окончен.
В противном случае, чтобы получить разложение в сумму простых фреймов Парсеваля, нужно применить
приведенный выше метод несколько раз. Введём оператор $D$ по следующей рекурсивной формуле:
$$
D(F,\,k) =
\left\{
\begin{array}{cl}
\lambda_{\alpha}  D(F_\alpha,\, k + 1) \oplus \lambda_{\beta}  D (F_\beta, k + 1), & F \text{ --- составной фрейм,} \\
F, & F \text{ --- простой фрейм.}
\end{array}\right.
$$

Остаётся доказать, что глубина рекурсии $k$ не превосходит некоторой константы. Согласно только что доказанному свойству
 при увеличении $k$ на единицу также на единицу уменьшается число векторов в фреймах $F_\alpha$ и $F_\beta$:
$$
|F_\alpha| \leq |F| - 1, \quad |F_\beta| \leq |F| - 1.
$$
Исходя из того, что фреймы Парсеваля с наименьшим числом векторов --- это ортонормированные базисы, получаем оценку
$k \leq M - N$, где $M$ "--- размерность фрейма при $k = 0$ , а $N$ "--- размерность пространства.
\end{proof}


\smallskip
\Section[n]{Равноугольные жёсткие фреймы}

\begin{definition}
Будем называть нормированный фрейм $F = \left\{f_i\right\}_{i=1}^M$ {\it равноугольным}, если существует константа
$c \in \left[\:0,1\right)$ такая что, для всех $i \leq j$ выполняется следующее равенство:
$$
\langle f_i, f_j\rangle =
\left\{
\begin{array}{rl}
1, &i = j, \\
\pm  c, & i \neq j.
\end{array}
\right.
$$
\end{definition}

В работе \cite{ryabtsov:Equ} доказывается, что равноугольная система является жёстким фреймом тогда и только тогда, когда
$$
c = \sqrt{\frac{M - N}{N(M - 1)}}.
$$

Для любой пары $M$ и $N$ существует не более одного равноугольного жесткого фрейма, причём для большинства пар таких фреймов нет.
Известными примерами равноугольных жёстких фреймов являются ортонормированные базисы, системы Мерседес"--~Бенц и другие \cite{ryabtsov:Equ}.

\begin{theorem}
Для любого равноугольного жесткого фрейма $F = \{f_i\}_{i=1}^M$ существует один и только один фрейм Парсеваля$,$
получаемый перенормировкой $\Bigl\{\sqrt{\frac{N}{M}} f_i \Bigr\}_{i=1}^M$ векторов фрейма$,$ причём этот фрейм простой.
\end{theorem}

\begin{proof}
Докажем, что система $\Bigl\{\sqrt{\frac{N}{M}} f_i \Bigr\}_{i=1}^M$ является фреймом Парсеваля.
Это утверждение является тривиальным следствием того, что граница любого нормированного жёсткого фрейма равна ${M}/{N}$ \cite{ryabtsov:Cas}.

Докажем, что полученный фрейм "--- простой. Для этого воспользуемся определением \eqref{ryab:eq2}.
Подставим вместо $x$ каждый из векторов перенормированного фрейма и получим необходимое
условие того, что фрейм является фреймом Парсеваля:
$$
k = 1, 2, \dots, M, \quad
\sum\limits_{i=1}^M \left|\Bigl\langle \sqrt{\tfrac{N}{M}} f_k, \alpha_i \sqrt{\tfrac{N}{M}} f_i\Bigr \rangle\right|^2 = 1
\Rightarrow 
\frac{N}{M}\sum\limits_{i=1}^M |\langle f_k, f_i \rangle|^2 |\alpha_i|^2 = 1.
$$

Введём обозначение $y_i = |\alpha_i|^2$. Получаем систему уравнений
$$
Cy = e,  \quad C = ({N}/{M}) \text{circ} \left\{1, c^2, \ldots, c^2\right\}, \quad
e =
\begin{pmatrix}
1, \, \ldots, \, 1 \\
\end{pmatrix}^{\top}.
$$

Матрица $C$ является невырожденной при $c^2 \neq 1$, система уравнений имеет единственное решение.
Подстановкой легко проверить, что решение всегда положительно и имеет следующий вид:
$$
y_k = (N / M) \left(c^2 (M - 1) + 1 \right), \quad
\alpha_k = \sqrt{ (N  /  M) \left(c^2 (M - 1) + 1 \right)}.
$$
Это противоречит определению составного фрейма, где требуется хотя бы одно значение $\alpha_k$, равное нулю, так как
это автоматически влечёт равенство нулю всех констант. Это противоречие доказывает, что полученный фрейм Парсеваля "--- простой.
\end{proof}

\begin{thebibliography}{5}

\Bibitem{ryabtsov:Cri}
\by Christensen~O.
\book An introduction to frames and Riesz bases. Applied and Numerical Harmonic Analysis
\publ Birkh\"auser Boston, Inc.
\publaddr Boston, MA
\yr 2003
\totalpages 440


\Bibitem{ryabtsov:Cas}
\by Casazza~P.\,G., Tremain~J.\,C.
\preprint A brief introduction to Hilbert-space frame theory and its applications
\preprintinfo preprint posted on \href{http://www.framerc.org/papers/IntroFrameTheory.pdf}{www.framerc.org}



\RBibitem{ryabtsov:Ist}
\by Истомина~М.\,Н., Певный~А.\,Б.
\paper О расположении точек на сфере и~фрейме Мерседес--Бенц
\serial Матем. просв., сер.~3
\yr 2007
\vol 11
\pages 105--112
\publ Изд-во МЦНМО
\publaddr М.
\misctransl
\by Istomina~M.\,N., Pevnyi~A.\,B.
\paper On disposition of points on a sphere and Mercedes--Benz frame
\serial Mat. Pros., Ser.~3
\yr 2007
\vol 11
\pages 105--112
\publ Izd-vo MCNMO
\publaddr Moscow


\Bibitem{ryabtsov:Me}
\by Novikov~S.\,Ya., Ryabtsov~I.\,S.
\paper Optimization of Frame Representations for Compressed Sensing and Mercedes-Benz Frame
\jour Proc. Steklov Inst. Math.
\yr 2009
\vol 265
\pages 199--207


\Bibitem{ryabtsov:Equ}
\by Casazza~P.~G., Redmond~D.,  Tremain~J.\,C.
\paper Real equiangular frames
\inbook Proc. $42$th Annu. Conf. Information Sciences and Systems $($CISS $2008)$
\publaddr Princeton, NJ
\yr 2008
\pages 715–720



\end{thebibliography}


\noindent
\hskip6.5cm \thedate \par\noindent
\hskip6.5cm\therevdate\par



\makeengtitle

\newpage


%\end{document}
